% Peekaboo-VL report: appendix_a_negative. Included from ../peekaboo_report.tex.
\section{What Did Not Work}
\label{sec:negative}
Negative results are useful for anyone working on this benchmark, so we
list them in one place.

\begin{itemize}
\item \textbf{Image-space wording of the motion hint}. The
measurements were right, but the words matched the wrong options. It cost
12.5 points on camera-motion questions.
\item \textbf{More frames or shifted frames at test time}. These
changes break the time layout the model learned in training, and made
exact-time questions about 5 points worse.
\item \textbf{Fixing pan detection with better thresholds}.
Pans and forward flight overlap too much in our simple four-number motion
model.
\item \textbf{Mixing in the text prior after fine-tuning}. It added
only 0.4 points, which could be luck.
\item \textbf{Letter debiasing after fine-tuning.} Because we shuffle
options in training, there is no letter bias left to remove; PriDe changes
nothing.
\item \textbf{Averaging over option orders at test time.} After
fine-tuning, showing the options in rotated orders and averaging the
answers gave no gain.
\item \textbf{Combining the fine-tuned model with models that were not fine-tuned.} An
ensemble with zero-shot models did not beat the fine-tuned model alone.
\item \textbf{A bigger model without fine-tuning.} A 4-bit Qwen3.8-27B scored
77.5\% on dev without fine-tuning, below the 9B model (79.6\%). Bigger is not
automatically better here, and the 27B model is too large to fine-tune on
our GPU.
\item \textbf{Switching to Cosmos-Reason2-8B.} With the same fine-tuning
recipe it ends 0.7 points below Qwen3.5-9B (test-weighted), mostly because
of Change Detection (Sec.~\ref{sec:backbone}).
\end{itemize}
